% ---------------------------------------------------------------------------------------------
% Supplementary theory (SI). Statements moved from the main-text theory section; notation as in sections/theory.tex.
% References to main-text labels use \Mref; proofs are in sections/appendix_proofs.tex (label si:proofs).
% ---------------------------------------------------------------------------------------------
\section{Additional theoretical results}\label{si:theory-extra}

This section collects the theoretical results that are not needed to follow the main text: the residual influence values, the
optimal coefficient and calibrated form of the difference estimator, closed-form design effects, misspecification bounds,
tuning of the difference coefficient, the zero-variance floor, allocation for the difference estimator, and the design effect as
a function of the disagreement rate. Equations and statements of the main text are cited by their main-text numbers.

\subsection{Residual influence values}\label{si:theory-resid}

The residual influence values of the difference estimator have an equally explicit form, and a
feature that the plain values lack: the label enters only through $Y_i-\pi_i$.

\begin{proposition}[Residual influence values]\label{prop:residual}
Let $d_i=\nabla f(T)^\top(u_i-v_i)$, which does not depend on the labels. Then $e_i=(Y_i-\pi_i)\,d_i$, and for $G_i=g$
\begin{align}
e^{TPR}_i&=s_g\frac{(Y_i-\pi_i)(D_i-\delta_g)}{P_g},\notag\\
e^{FPR}_i&=s_g\frac{\{(1-Y_i)-(1-\pi_i)\}(D_i-\gamma_g)}{Q_g},\notag\\
e^{ERR}_i&=-\frac{2(Y_i-\pi_i)D_i}{N},\label{eq:resid-forms}
\end{align}
and $\tilde z_i=z_i-e_i$ equals (\Mref{eq:infl-forms}) with $Y_i$ replaced by $\pi_i$. In a stratum where the pair agrees,
$e^{ERR}_i=0$ for every $i$: the stratum contributes neither variance nor error to the difference estimator of $\theta^{ERR}$.
Under the working model that $Y_1,\dots,Y_N$ are independent with $P(Y_i=1)=\pi_i$ (the proxy is calibrated), with expectation $E_\pi$,
the gradient evaluated at the non-random model-expected totals $T_\pi=\sum_{i\in U}\psi_i$ (as in the design; this changes $d_i$ by a relative $O(r)$ only),
and with $\bar\sigma^2_h=N_h^{-1}\sum_{U_h}\pi_i(1-\pi_i)d_i^2$ and $C_h=\mathrm{Cov}_h(z,\tilde z)$ the stratum covariance,
\begin{align}
&E_\pi(e_i)=0,\quad \mathrm{Var}_\pi(e_i)=\pi_i(1-\pi_i)d_i^2,\notag\\
&E_\pi S^2_{e,h}=\bar\sigma^2_h,\quad E_\pi S^2_{z,h}=S^2_{\tilde z,h}+\bar\sigma^2_h,\notag\\
&E_\pi C_h=S^2_{\tilde z,h}.\label{eq:resid-moments}
\end{align}
\end{proposition}

\subsection{Difference estimator: coefficient and calibrated form}\label{si:theory-dif}

Recall the family $\hat T_{\lambda\psi}$ obtained by replacing $\psi$ by $\lambda_h\psi$ in stratum $h$ ($\lambda\equiv0$ is the plain estimator and $\lambda\equiv1$
the difference estimator); every member is unbiased for $T$ by Proposition~\Mref{prop:validity}(a), and $\beta_h$, $R_h$ are as in Section~\Mref{sec:theory-dif}.
The sufficient condition of Proposition~\Mref{prop:dif-var} can be restated as $R_h>S_{\tilde z,h}/(2S_{z,h})$ in every stratum: the within-stratum correlation of $z_i$ with its proxy mean must exceed half the ratio of their standard deviations.

\begin{proposition}[Variance of the difference estimator, continued]\label{prop:dif-var-ext}
Under the notation of Proposition~\Mref{prop:dif-var}, for any label-free proxy and any allocation:
\begin{enumerate}
\item[(a)] \emph{Coefficient.} $V(\lambda)=\sum_hc_h\{S^2_{z,h}(1-R_h^2)+(\lambda_h-\beta_h)^2S^2_{\tilde z,h}\}$. It is minimised at
$\lambda_h=\beta_h$, with $V(\beta)=\sum_hc_hS^2_{z,h}(1-R_h^2)$; moreover
$V_0=V(\beta)+\sum_hc_h\beta_h^2S^2_{\tilde z,h}$ and $V_\psi=V(\beta)+\sum_hc_h(1-\beta_h)^2S^2_{\tilde z,h}$. With the clipped coefficient
$\lambda_h=\min\{1,\max(0,\beta_h)\}$, $V(\lambda)\le\min(V_0,V_\psi)$.
\item[(b)] \emph{Calibrated proxy, design-based form.} If $\mathrm{Cov}_h(e,\tilde z)=0$ in every stratum (the residual is uncorrelated with the
proxy influence within strata), then $\beta_h=1$ and $V_0-V_\psi=\sum_hc_hS^2_{\tilde z,h}\ge0$ exactly.
\end{enumerate}
In terms of correlations, $V_0-V_\psi=\sum_hc_h\,S_{\tilde z,h}\,(2R_hS_{z,h}-S_{\tilde z,h})$.
\end{proposition}

Part (a) of Proposition~\Mref{prop:dif-var} shows that, unlike what is often assumed for prediction-powered estimators, the sign of the gain is not automatic:
$\lambda=1$ is the right coefficient only if the proxy influence is neither over- nor under-dispersed relative to
its regression on $z$; the harm from an over-dispersed proxy is quantified in Proposition~\ref{prop:dif-proxy}.

\subsection{Design-effect ratio: closed forms and discussion}\label{si:theory-deff-ratio}

The following completes Proposition~\Mref{prop:dif-deff}; its proof is in Supplementary Section~\ref{si:proofs}.

\begin{proposition}[Closed forms of the stratum gain]\label{prop:dif-closed}
Under the setting of Proposition~\Mref{prop:dif-deff}, for $\theta^{TPR},\theta^{FPR}$ and $\theta^{ERR}$ the stratum gains depend on the proxy as follows. Let $\overline{\pi(1-\pi)}$ denote
the stratum mean of $\pi_i(1-\pi_i)$. In a \emph{disagreement stratum} ($|D_i|=1$ on $U_h$, $D_i$ and $\pi_i$ uncorrelated within the stratum,
$\bar a$ the stratum mean of $D_i$, and $\delta_g,\gamma_g$ negligible),
\begin{align}
1+\nu^{TPR}&=\frac{\bar\pi(1-\bar\pi\bar a^2)}{\overline{\pi(1-\pi)}},\notag\\
1+\nu^{FPR}&=\frac{(1-\bar\pi)\{1-(1-\bar\pi)\bar a^2\}}{\overline{\pi(1-\pi)}},\notag\\
1+\nu^{ERR}&=\frac{1-\bar a^2(1-2\bar\pi)^2}{4\,\overline{\pi(1-\pi)}} .\label{eq:dif-closed}
\end{align}
In an \emph{agreement stratum} ($D_i=0$) the error difference has $e_i=\tilde z_i=0$, and for the rate differences
$1+\nu=\{\overline{\pi(1-\pi)}+S^2_{\pi,h}\}/\overline{\pi(1-\pi)}$, the inverse of the share of the Bernoulli variance that the proxy does \emph{not} explain.
\end{proposition}

Four consequences. First, calibration alone guarantees $G\ge1$: with a calibrated proxy the difference estimator is never worse than the plain one in
expectation, for any allocation. Second, the gain is a within-stratum $R^2$, so it is governed by what the proxy influence $\tilde z_i$ explains of $z_i$
\emph{inside} a stratum, and the strata are already built to remove between-stratum variation. Third, $\tilde z_i$ contains two kinds of
information: the proxy $\pi_i$ and the known structure of $u_i,v_i$ (the signs and the pattern of $D_i$). The second kind is free and
is the main source of the gain: in a two-sided disagreement stratum with a constant proxy $\pi_i\equiv\bar\pi$ and $\bar a=0$,
\eqref{eq:dif-closed} gives gains $1/(1-\bar\pi)$ for $\theta^{TPR}$, $1/\bar\pi$ for $\theta^{FPR}$ and $1/\{4\bar\pi(1-\bar\pi)\}$ for $\theta^{ERR}$, i.e.\ 1.4,
3.3 and 1.2 for $\bar\pi=0.3$ and exactly 2, 2 and 1 for $\bar\pi=1/2$. Only the error difference needs a sharp proxy: its gain is
$\{1-\overline{(1-2\pi)^2}\}^{-1}$ and equals one for $\pi\equiv1/2$, which makes $\tilde z\equiv0$ and the two estimators identical. Fourth, the
gain is bounded for proxies of realistic sharpness, because the disagreement strata are the strata in which the compared models, and hence typically the
averaged probabilities $\pi_i$, are least certain: a large $G$ requires within-stratum $R_h^2$ close to one, that is a sharp proxy
($\pi_i$ near 0 or 1) and a stratum in which the models still disagree. Where the proxy is that sharp the predictions would be nearly
exact and few labels would be needed in any case. For disagreement strata with $\bar\pi$ between 0.2 and 0.5 the closed forms give gains of the order of one to three, on top of the
stratification gain; Section~\ref{sec:sim-dif} compares the prediction with the Monte Carlo ratios.

Gain and stratification multiply: for the same sample the efficiency of the difference estimator relative to SRS is, to the order of the
approximations of Section~\Mref{sec:theory-deff}, $\mathrm{Eff}_t\cdot G_t$, with $\mathrm{Eff}_t$ from
(\Mref{eq:deff-general}) and $G_t$ from (\Mref{eq:dif-deff}). They are not independent: a stratifier that
removes more of the proxy's within-stratum signal (a finer partition of the proxy range) reduces $G$ and raises $\mathrm{Eff}_t$.

\subsection{Misspecified proxy for the difference estimator}\label{si:theory-dif-proxy}

\begin{proposition}[Misspecified proxy]\label{prop:dif-proxy}
\begin{enumerate}
\item[(a)] For every proxy, $\hat T_\psi$ is design-unbiased and $\hat V_\psi$ is a consistent variance estimator
(Proposition~\Mref{prop:validity}). Misspecification can affect efficiency only.
\item[(b)] Let the labels be independent with $P(Y_i=1)=q_i$, where $q_i$ may differ from $\pi_i$, and let $\tilde z^{q}_i$, $\tilde z^{\pi}_i$ be the proxy
influence values for $q$ and $\pi$. With $E_q$ the expectation under this model,
\begin{equation}\label{eq:dif-misspec}
E_q(V_0-V_\psi)=\sum_hc_h\Big\{S^2_{\tilde z^{q},h}-S^2_{\tilde z^{q}-\tilde z^{\pi},h}\Big\},
\end{equation}
so the difference estimator wins if and only if, within strata and on the $c_h$-weighted scale, the proxy error in influence units is smaller than the signal it
should capture. If $\tilde z^{\pi}_i-\bar{\tilde z}^{\pi}_h=\kappa(\tilde z^{q}_i-\bar{\tilde z}^{q}_h)$ in every stratum (a proxy that is sharper by the factor
$\kappa$ than the truth), the estimator wins if and only if $0<\kappa<2$; for $\kappa<0$ (anti-correlated proxy) it loses.
\item[(c)] In all cases $V_\psi\le(1+\max_hS_{\tilde z,h}/S_{z,h})^2\,V_0$, with equality only if $\tilde z_i=-s\,z_i$ within each stratum
for a constant $s>0$ (perfect anti-correlation). Hence the possible loss is bounded by a function of the dispersion of the proxy influence relative to the plain one.
\item[(d)] With the clipped coefficient of Proposition~\ref{prop:dif-var-ext}(a) the variance never exceeds that of the plain estimator
(to first order). In particular, an uninformative proxy, in the sense $\beta_h=0$ ($\tilde z$ uncorrelated with $z$ within strata), has optimal coefficient $\lambda_h=0$ and
then $V_\psi=V_0+\sum_hc_hS^2_{\tilde z,h}$: the fixed coefficient $\lambda=1$ adds the proxy's own variance.
\end{enumerate}
\end{proposition}

Two clarifications. The statement that the difference estimator can be worse than the plain estimator \emph{only} if the proxy is anti-correlated is
too strong: by (\Mref{eq:dif-identity}) it is worse whenever $\beta_h<1/2$, so a weakly aligned or over-dispersed
proxy ($\kappa>2$) also hurts, while a modestly noisy but correctly scaled proxy does not. And ``uninformative'' must be read in terms of $\tilde z_i$,
not of $\pi_i$: a constant proxy $\pi_i\equiv\bar\pi$ still has $\tilde z_i$ varying with the known pattern of $D_i$, and for the rate differences
it still gains, as the closed forms \eqref{eq:dif-closed} show; only for $\theta^{ERR}$ with $\bar\pi=1/2$ is $\tilde z\equiv0$, and
only an inverted proxy ($\kappa<0$) makes $\beta_h<0$ for the error difference. Section~\ref{sec:sim-dif} confirms both points numerically.

\subsection{Fixed versus tuned coefficient}\label{si:theory-lambda}

\paragraph{Fixed versus tuned coefficient.} In prediction-powered inference \citep{angelopoulos2026ppipp,fisch2024stratppi} and in the
regression-estimator literature \citep{cassel1976,sarndal1992} the coefficient is estimated from the labelled sample, here
per stratum and per estimand: $\hat\lambda_h=\widehat{\mathrm{Cov}}(\hat z,\hat{\tilde z})/\widehat{\mathrm{Var}}(\hat{\tilde z})$ clipped to $[0,1]$.

\begin{proposition}[Estimated coefficient]\label{prop:dif-lambda}
Suppose $\hat\lambda_h$ is a function of the labelled sample of stratum $h$ with $\hat\lambda_h-\lambda^\infty_h\to0$ in probability, for deterministic
$\lambda^\infty_h$ (for the clipped sample slope, $\lambda^\infty_h=\min\{1,\max(0,\beta_h)\}$ whenever $\beta_h$ converges). Under (A1)--(A4)
for the fixed coefficient $\lambda^\infty$, $\hat\theta_{\hat\lambda\psi}-\hat\theta_{\lambda^\infty\psi}=o_p(V^{1/2})$; hence the tuned estimator has the same asymptotic
normal limit and linearised variance $V(\lambda^\infty)$ as the one with the optimal fixed coefficient, and is never asymptotically worse than the plain
estimator. At second order, estimating $\hat\lambda_h$ adds to the variance of stratum $h$ approximately
$c_h(1-R_h^2)S^2_{z,h}/n_h$, a relative increase of order $1/n_h$ in a stratum that is already optimally tuned, and the plug-in estimator $\hat V_\psi$ computed
with the fitted $\hat\lambda_h$ does not account for it (it is biased downwards by a factor of the same order).
\end{proposition}

The consequence is a trade-off that is independent of the particular data. The best that tuning can add to $\lambda=1$ is the factor
$V_\psi/V(\lambda^\ast)=1+\sum_hc_h(1-\beta_h)^2S^2_{\tilde z,h}/V(\beta)$ (\Mref{eq:dif-identity}), which is small when the proxy is nearly calibrated ($\beta_h\approx1$): then
tuning gains little in variance but costs an $O(1/n_h)$ inflation in each of the $H$ strata and a downward-biased variance estimator, which matters when $n_h$ is small, as
at the budgets of this paper ($n_h$ of the order of 10). With a miscalibrated proxy the gain from tuning can be real (see the $\kappa\ne1$ cases in
Section~\ref{sec:sim-dif}), but it is paid for in coverage. The calibrated difference estimator with $\lambda=1$ has no estimated
parameters, exact unbiasedness and a plug-in variance that is unbiased for the totals (Proposition~\Mref{prop:validity}(a)), which is why we adopt it as the main
estimator. Section~\ref{sec:sim-dif} reports the predicted and the observed gain of tuning together with the coverage of each version.

\subsection{Zero-variance strata and the variance floor}\label{si:theory-floor}

The main text states the floor briefly (Remark~\Mref{rem:zero-var}); the full discussion follows.

\begin{remark}[Zero-variance strata and the variance floor, detailed]\label{rem:zero-var-ext}
The Wald interval needs $\hat V_\psi>0$ and $\hat V_\psi\approx V_\psi$, and both fail when a stratum's sample shows no variation in $\hat e_i$ although
$S^2_{e,h}>0$. For the difference estimator this has three sources. (i) \emph{Structural zeros}: strata with $d_i\equiv0$ (the error
difference where the pair agrees, Proposition~\ref{prop:residual}), where $e_i\equiv0$ in the whole stratum; the estimate is exact there and
$\hat V_\psi$ correctly contributes zero. (ii) \emph{Sampling zeros}: a stratum with few positives (or negatives) in which the sample has none. For the plain estimator $\hat z_i=0$ for every sampled unit, so $s^2_{\hat z,h}=0$ exactly;
for the difference estimator $\hat e_i=-\pi_i\hat d_i$ varies with $\pi_i\hat d_i$, so exact zeros are rarer, but the sample variance is still far below the
population variance $\bar\sigma^2_h=N_h^{-1}\sum_{U_h}\pi_i(1-\pi_i)d_i^2$ whenever the missed positives matter. With a calibrated proxy of level $\pi$ and constant $d$, the
probability of a sample without positives is $(1-\pi)^{n_h}$, which is the zero-count failure of the Wald interval for a binomial proportion \citep{brown2001}.
(iii) \emph{Exact proxies}: if $\pi_i\in\{0,1\}$ equals $Y_i$ then $e_i\equiv0$ and
(A3) fails; the estimator is then exact, but the proxy is unlikely to be exact on the whole pool. The floor of Section~\Mref{sec:ablations} carries over: replace $s^2_{\hat e,h}$ by the anticipated
residual variance
\begin{align}
\tilde s^2_{h}&=N_h^{-1}\sum_{i\in U_h}\pi_i(1-\pi_i)\hat d_i^{\,2},\label{eq:floor}\\
\hat d_i&=\nabla f(\hat T_\psi)^\top(u_i-v_i),\notag
\end{align}
whenever $s^2_{\hat e,h}=0$ (or, more conservatively, use $\max\{s^2_{\hat e,h},\eta\tilde s^2_h\}$ with $\eta\in(0,1]$). Because $\tilde s^2_h=0$ exactly when $d_i\equiv0$ on $U_h$, structural zeros
are untouched. The floor is label-free, never decreases $\hat V_\psi$, leaves the point estimate unchanged, and (Supplementary Section~\ref{si:proofs}) does not affect
Proposition~\Mref{prop:validity}(c)--(d) under Assumption~(A5) (the floor is asymptotically negligible): the floor is active with probability tending to zero in every stratum with $S^2_{e,h}$ bounded away from zero, and in strata
with vanishing variance it is, if anything, conservative. It is a finite-sample safeguard, not an asymptotic correction, and Section~\ref{sec:sim-dif} reports its effect on coverage.
\end{remark}

\subsection{Allocation for the difference estimator}\label{si:theory-alloc}


In the design the $S^2_{ht}$ are unknown. They are replaced by their expectations under the
working model $Y_i\sim\text{Bernoulli}(\pi_i)$ independently, with the gradient evaluated at the
model-expected totals; up to terms of order $1/N_h$ this is the anticipated variance of
\citet{isaki1982}. By \eqref{eq:resid-moments} the anticipated stratum variances are (\Mref{eq:anticipated}), with
$\bar\sigma^2_{ht}=N_h^{-1}\sum_{U_h}\pi_i(1-\pi_i)d_{it}^2$. The plain rule is what the design uses (Section~\Mref{sec:design}). The residual rule is the Neyman allocation for the prediction-powered residual
\citep{fisch2024stratppi}; it depends on the proxy only through $\pi_i(1-\pi_i)$ and the label-free $d_i$, so a proxy of any
sharpness gives a well-defined allocation.

\begin{proposition}[Allocation for the difference estimator]\label{prop:alloc-dif}
Let $n^\psi$ be the optimal allocation for the difference estimator ($A_h=A^\psi_h$, $A^{\psi\,2}_h=\sum_tw_tS^2_{e,ht}$) and $\tilde n^0$ the allocation obtained from the
plain anticipated variances. Put $\tilde\nu_h=\sum_tw_tS^2_{\tilde z,ht}/\sum_tw_t\bar\sigma^2_{ht}$, $x_h=(1+\tilde\nu_h)^{1/2}$ and assume the anticipated values are exact for the
residual variances ($\tilde A^\psi_h=A^\psi_h$). Ignoring caps, minima and rounding, the weighted variance of the difference estimator with the \emph{wrong} (plain) allocation satisfies
\begin{align}
\frac{\Phi_\psi(\tilde n^0)}{\Phi_\psi(n^\psi)}&=\Big(\sum_h\omega_hx_h\Big)\Big(\sum_h\omega_h/x_h\Big)\notag\\
&\le\frac{(x_{\max}+x_{\min})^2}{4x_{\max}x_{\min}},\notag\\
\omega_h&=\frac{N_hA^\psi_h}{\sum_kN_kA^\psi_k},\label{eq:alloc-loss}
\end{align}
with equality to one iff $\tilde\nu_h$ is constant over the strata with $\omega_h>0$. The same expressions hold, with $x_h$ replaced by $1/x_h$ and
$\omega_h\propto N_hA^0_h$, for the plain estimator allocated with the residual variances. Consequently, if $\tilde\nu_h\in[0,\nu_{\max}]$, the loss from
using the wrong rule is at most $\{1+(1+\nu_{\max})^{1/2}\}^2/\{4(1+\nu_{\max})^{1/2}\}$ in either direction, for instance $9/8$ for $\nu_{\max}=3$.
\end{proposition}

So the choice between the two rules is a second-order matter wherever the proxy's signal-to-noise ratio $\tilde\nu_h$ is similar across strata, which is the case for strata with
comparable disagreement structure, and it matters when $\tilde\nu_h$ is large in some strata and small in others (for example, very sharp proxies in the agreement strata of the
rate differences). Neither rule is optimal for the other estimator, and the plain rule remains the sensible default for the difference estimator: it is exactly Neyman for the plain estimator, is
robust to a proxy whose calibration is not trusted (the residual rule depends on $\pi_i(1-\pi_i)$ and gives zero weight to strata in which the proxy is
extreme), and its efficiency loss for the difference estimator is bounded by \eqref{eq:alloc-loss}. Section~\ref{sec:sim-dif} reports the predicted and the observed loss.

\subsection{Design effect and the disagreement rate}\label{si:theory-deff}

Let $\sigma^2_{ht}=(1-1/N_h)S^2_{ht}$ and $\sigma^2_t=(1-1/N)S^2_t$ be the variances with divisor
$N_h$ and $N$. For a single target allocated by its own Neyman rule, and with $L/N\to0$,
the efficiency relative to SRS is (\Mref{eq:deff-general}) of the main text,
$\mathrm{Eff}_t=\sigma_t^2/(\sum_hW_h\sigma_{ht})^2\ge1$, because
$\sigma_t^2=\sum_hW_h\sigma_{ht}^2+\sum_hW_h(\bar z_h-\bar z)^2\ge(\sum_hW_h\sigma_{ht})^2$; equality
holds only if all within-stratum standard deviations are equal and the stratum means coincide.

Combined with Proposition~\Mref{prop:influence}, equation~(\Mref{eq:deff-general}) gives explicit formulas for the plain estimator.
We state them for an idealised design in which, within each group, the strata separate the
instances on which $m$ and $j$ disagree (set $\mathcal D_g$, share $\rho_g$ of group $g$) from
those on which they agree (set $\mathcal A_g$). This is exactly the pattern design when $M=2$; for
$M>2$ it describes a design stratified for the pair $(m,j)$, and is discussed below.
Finite-population corrections are ignored ($L/N\to0$ and $n_h/N_h\to0$); they are retained in
the exact linearised efficiency used in Section~\ref{sec:simulation}.

\begin{corollary}[Design effect as a function of the disagreement rate]\label{prop:deff}
Let $\zeta_g=N_g/N$, $p_g=P_g/N_g$, $r_g=N_g^{-1}\sum_{G_i=g}Y_i|D_i|$ (positives on which the pair
disagrees, as a share of the group), $\pi_{\mathcal D,g}=r_g/\rho_g$,
$\pi_{\mathcal A,g}=(p_g-r_g)/(1-\rho_g)$ and $a_g=c_{\mathcal D,g}/\pi_{\mathcal D,g}$, where
$c_{\mathcal D,g}=\delta_gp_g/\rho_g$ is the mean of $Y_iD_i$ over $\mathcal D_g$; $a_g\in[-1,1]$
is the net share of disputed positives on which $m$ is right and $j$ wrong.
\begin{enumerate}
\item[(a)] For $\theta^{TPR}_{mj}$, with own-target Neyman allocation on the strata
$\{\mathcal A_g,\mathcal D_g\}_{g=0,1}$ and $L/N\to0$,
\begin{equation}\label{eq:deff-tpr}
\mathrm{Eff}^{TPR}=\frac{\sum_g(r_g-\delta_g^2p_g)/(\zeta_gp_g^2)}
{\big[\sum_g\{\rho_g\sigma_{\mathcal D,g}+(1-\rho_g)\sigma_{\mathcal A,g}\}/p_g\big]^2},
\end{equation}
$\sigma_{\mathcal A,g}=|\delta_g|\{\pi_{\mathcal A,g}(1-\pi_{\mathcal A,g})\}^{1/2}$ and $\sigma^2_{\mathcal D,g}=\pi_{\mathcal D,g}-c_{\mathcal D,g}^2+\delta_g^2\pi_{\mathcal D,g}(1-\pi_{\mathcal D,g})
-2\delta_gc_{\mathcal D,g}(1-\pi_{\mathcal D,g})$. The same holds for $\theta^{FPR}_{mj}$ with $Y$
replaced by $1-Y$. For $\theta^{ERR}_{mj}$, with $\rho=\sum_g\zeta_g\rho_g$ and $\mu_g$ the mean of
$D_i(1-2Y_i)$ over $\mathcal D_g$,
\begin{equation}\label{eq:deff-err}
\mathrm{Eff}^{ERR}=\frac{\rho-(\theta^{ERR}_{mj})^2}{\big(\sum_g\zeta_g\rho_g\sqrt{1-\mu_g^2}\big)^2}.
\end{equation}
Any further split of $\mathcal A_g$ or $\mathcal D_g$ (for example $\mathcal A_g$ into $K=0$ and
$K=M$) can only increase \eqref{eq:deff-tpr} and \eqref{eq:deff-err}.
\item[(b)] \emph{Small targets.} If $\delta_g=0$ and $\rho_g=\rho$ for both groups, then
\begin{equation}\label{eq:deff-small}
\mathrm{Eff}^{TPR}=\frac1\rho\,\Gamma,\qquad
\Gamma=\frac{\sum_gr_g/(\zeta_gp_g^2)}{(\sum_g\sqrt{r_g}/p_g)^2}\ge1,
\end{equation}
where $\Gamma$ is the efficiency of Neyman allocation between the two groups alone; and if
$\theta^{ERR}_{mj}=0=\mu_g$, then $\mathrm{Eff}^{ERR}=1/\rho$.
\item[(c)] \emph{Target size.} For one group (drop $g$), with $\pi_{\mathcal D}$,
$\pi_{\mathcal A}$ and $a$ held fixed, as $\rho\to0$,
\begin{align}
\mathrm{Eff}^{TPR}&=\frac{1+O(\rho)}{\rho\,(\beta+\lambda)^2},\notag\\
\beta&=\sqrt{1-a^2\pi_{\mathcal D}},\notag\\
\lambda&=\frac{|a|\sqrt{\pi_{\mathcal D}\,\pi_{\mathcal A}(1-\pi_{\mathcal A})}}{p}.\label{eq:deff-lambda}
\end{align}
Since $\delta=\rho\pi_{\mathcal D}a/p$, the target is proportional to $\rho$ and to $a$. A more
one-sided disagreement set (larger $|a|$, hence a larger target) changes the gain by the bounded
factor $(\beta+\lambda)^{-2}$, which is below one when the agreement region is noisy
($\lambda>1-\beta$), but it does not change the order $1/\rho$.
\item[(d)] \emph{When the gain vanishes.} As $\rho_g\to1$ the agreement strata disappear and
\eqref{eq:deff-tpr} and \eqref{eq:deff-err} tend to the efficiency of stratification by
group alone (for the rate differences with $\delta_g=0$ this is $\Gamma$; for the error difference
it equals 1 when $\mu_0=\mu_1$). The gain beyond group stratification is therefore governed by $1-\rho$: it is large when
the compared models disagree on a small share of the pool, it is reduced by a constant factor
when the pair's disagreements are one-sided ($|a_g|$ near 1) and the agreement region has
substantial label noise, and it disappears when the models disagree almost everywhere.
\end{enumerate}
\end{corollary}

The heuristic \emph{efficiency $\approx1/\rho$} is thus exact for an error difference with a
two-sided disagreement set ($\mu_g=0$), and correct up to the group factor $\Gamma$ and the
target-size factor $(\beta+\lambda)^{-2}$ for the rate differences. For a pair that disagrees on
5\% of the pool it promises a gain of the order of 20 relative to SRS, and of 2.5 for a pair that
disagrees on 40\%. These are gains of the stratification; the difference estimator multiplies them by the factor $G$ of
(\Mref{eq:dif-deff}), which by \eqref{eq:dif-closed} is of order one to three and does not grow as $\rho\to0$: the difference estimator removes the part of the
label noise that a proxy of moderate sharpness can explain, whereas the stratification removes the unlabelled agreement region.

For $M>2$ the pattern strata $(G_i,K_i)$ do not separate the disagreements of a particular pair:
an instance with $K_i=4$ out of nine may or may not be disputed by $m$ and $j$. The general
formula (\Mref{eq:deff-general}) still applies with the stratum variances of
Proposition~\Mref{prop:influence}. For an error difference whose disagreements are two-sided within
every stratum ($e_h=0$ for all $h$, hence $\theta^{ERR}_{mj}=0$), $\sigma^2_{ht}\propto\rho_h$
and (\Mref{eq:deff-general}) becomes
\begin{equation}\label{eq:deff-mixed}
\mathrm{Eff}^{ERR}=\frac{\rho}{(\sum_hW_h\sqrt{\rho_h})^2},\qquad 1\le\mathrm{Eff}^{ERR}\le\frac1\rho,
\end{equation}
since $\rho_h\le\sqrt{\rho_h}$ and $(\sum_hW_h\sqrt{\rho_h})^2\le\sum_hW_h\rho_h=\rho$. The bound $1/\rho$ is
attained only when every stratum is either entirely disputed or entirely agreed; the pattern
strata approach it to the extent that the pair's disagreements concentrate in a few values of
$K$. A lower bound of the same type holds for the pattern strata. Let $\rho_{\rm any}$ be the share
of instances on which not all $M$ models agree, i.e.\ with $0<K_i<M$. The pair can disagree only
there, so $\rho_h=0$ in the strata with $K\in\{0,M\}$, and Cauchy--Schwarz over the remaining strata
gives $(\sum_hW_h\sqrt{\rho_h})^2\le\rho_{\rm any}\,\rho$, hence
\begin{equation}\label{eq:deff-sandwich}
\frac{1}{\rho_{\rm any}}\;\le\;\mathrm{Eff}^{ERR}\;\le\;\frac{1}{\rho},
\end{equation}
with equality on the left when the pair's disagreements are spread evenly over the strata with
$0<K<M$ ($\rho_h=\rho/\rho_{\rm any}$ there). With nine models $\rho_{\rm any}$ is large (between
0.18 and 0.73 in our pools) while the pairwise $\rho$ is small, and the position of a pair within
this interval measures how well the counts $K$ locate its disagreements. Two further losses separate \eqref{eq:deff-tpr} and \eqref{eq:deff-err} from the efficiency of the implemented
design: a single allocation serves all $3\binom M2$ comparison estimands
(Proposition~\Mref{prop:allocation}(d)), and it is computed from the proxy
(Section~\Mref{sec:theory-proxy}). The exact linearised efficiency of the implemented design is
$N^2(1/L-1/N)S^2_t/V_t(n)$ with $V_t$ from (\Mref{eq:linvar}) (or (\Mref{eq:resid}) for the difference estimator); Section~\ref{sec:simulation}
compares it, and the closed forms above, with Monte Carlo results.

\subsection{Imperfect proxy in the allocation}\label{si:theory-proxy}

Let $\tilde n$ be the allocation computed from the anticipated values $\tilde A_h$, and $A_h$ the
values computed from the true labels, both with the same weights. This subsection concerns the proxy in the allocation; its role in the
difference estimator is Proposition~\ref{prop:dif-proxy}.

\begin{proposition}[Imperfect proxy]\label{prop:proxy}
\begin{enumerate}
\item[(a)] Proposition~\Mref{prop:validity} holds for $\tilde n$ whatever the proxy, provided
$\tilde n_h\ge2$ in every stratum (guaranteed by the minimum of 10 labels).
\item[(b)] Ignoring caps, minima and rounding, the relative loss in the weighted variance is
\begin{align}
\frac{\Phi_0(\tilde n)}{\Phi_0(n^\ast)}&=
\frac{(\sum_hN_h\tilde A_h)(\sum_hN_hA_h^2/\tilde A_h)}{(\sum_hN_hA_h)^2}\notag\\
&=\Big(\sum_h\omega_hx_h\Big)\Big(\sum_h\omega_h/x_h\Big)\ge1,\label{eq:proxy-loss}
\end{align}
with $\omega_h=N_hA_h/\sum_kN_kA_k$ and $x_h=\tilde A_h/A_h$. The loss is zero if and only if
$\tilde A_h\propto A_h$; a proxy that is wrong by a common factor costs nothing.
\item[(c)] If $x_{\min}\le x_h\le x_{\max}$, the loss is at most
$(x_{\max}+x_{\min})^2/(4x_{\max}x_{\min})$ (Kantorovich inequality). For instance, if every
anticipated value $\tilde A_h$ is within a factor 2 of $A_h$ up to a common scale
($x_{\max}/x_{\min}\le4$), the weighted variance is at most $25/16$ times the optimum on the same
strata, so the design keeps at least 64\% of the efficiency of the optimal allocation, whatever
the number of strata.
\item[(d)] Proportional allocation is the special case $\tilde A_h\equiv1$. Neither allocation
dominates the other for every population: \eqref{eq:proxy-loss} is smaller for the proxy
allocation if and only if
\begin{align*}
&\Big(\sum_h\omega_hx_h\Big)\Big(\sum_h\omega_h/x_h\Big)\\
&\quad\le\Big(\sum_h\omega_h/A_h\Big)\Big(\sum_h\omega_hA_h\Big).
\end{align*}
A proxy that is badly wrong in a high-variance stratum can make the design worse than
proportional allocation, and, through the minimum of $\ell$ labels in that stratum, worse than SRS.
\item[(e)] A defensive mixture $n_h=\eta\,n^{\rm prop}_h+(1-\eta)\tilde n_h$, $0<\eta<1$, satisfies
$\Phi_0(n)\le\Phi_0(n^{\rm prop})/\eta$ and $\Phi_0(n)\le\Phi_0(\tilde n)/(1-\eta)$; since
proportional allocation is never worse than SRS up to the factor $(1-1/\min_hN_h)^{-1}$, the
mixture is never worse than $1/\eta$ times SRS in the weighted variance without finite-population corrections, and with them (sampling fraction $f=L/N$) never worse than $(1-\eta f)/\{\eta(1-f)\}$ times SRS,
both up to the same factor; the second bound is $\eta^{-1}\{1+O(f)\}$.
\end{enumerate}
\end{proposition}

What can be guaranteed is therefore validity (part (a)) and a loss that is controlled by how far
the anticipated stratum standard deviations are from proportional to the true ones (parts (b) and
(c)). What cannot be guaranteed for the implemented rule is improvement over SRS for every
population; that requires a defensive component as in part (e). Two features make large losses
unlikely in our setting. The dominant part of the structure, namely which strata carry the
numerators of the comparison estimands, is determined by the predictions and does not depend on the
proxy (Proposition~\Mref{prop:influence}); and the proxy enters $\tilde A_h$ only through
$\pi_h(1-\pi_h)$ and $r_h$, which are smooth in the proxy probabilities. Section~\ref{sec:simulation}
quantifies the loss for good, uninformative and adversarial proxies.